\section{Inverse dynamics modeling loss}
\label{app:ablations}

We add inverse dynamics to improve agent localization and planning, while
leaving cue binding unsupervised.

\subsection{Predictive and inverse-dynamics objectives}

We train randomly initialized ViT-tiny models on 10,000 successful
trajectories for ten epochs, with batch size 128 and seed 3072. For a
192-dimensional projected state $z_t$ and its prediction $\hat z_{t+1}$,
\[
\mathcal L=\mathcal L_{\mathrm{pred}}+.09\mathcal L_{\mathrm{SIGReg}}
+\lambda_{\mathrm{IDM}}\mathcal L_{\mathrm{IDM}},\qquad
\mathcal L_{\mathrm{pred}}=\frac{1}{192}\mathbb E\|\hat z_{t+1}-z_{t+1}\|_2^2,
\]
\[
\mathcal L_{\mathrm{IDM}}=\frac{1}{10}\mathbb E\|h([z_{t+1},z_t])-a_t\|_2^2.
\]
SIGReg encourages isotropic Gaussian states. The target $a_t$ contains five
normalized 2-D actions. The training-only inverse head is
Linear$(384,512)$, LayerNorm, GELU, and Linear$(512,10)$; its gradients pass
through both states. We use $\lambda_{\mathrm{IDM}}=3.125$, versus zero for
the predictive objective without IDM. The zero-weight controls retain the
same inverse head but disable its loss. Inverse dynamics and auxiliary-task theory motivate
preserving action-relevant features
\citep{ivashkov2026sensorimotor,yu2025auxjepa}; our head receives no cue or
portal-pairing labels.

\subsection{Agent localization and planning}

Adding IDM improves agent-position probes (Table~\ref{tab:agentidm}) and
planning (Table~\ref{tab:idmplanning}). With spatial aggregation, pooled
CEM/GD success rises from 8/8\% to 48/58\%; with CLS, it rises from 4/8\%
to 38/40\%. Within each readout, initialization, training budget, and
evaluation are matched, and only the IDM weight changes.

\begin{table}[!htbp]
\centering
\caption{\textbf{IDM makes the agent's position more accessible.}
Linear probes of agent coordinates on the same four Color checkpoints as
Table~\ref{tab:idmplanning}. Current and predicted states target the current
and next agent position, respectively; MAE is in pixels and $R^2$ averages
$x/y$ scores. Probes fit on 640 episodes and evaluate on 256 disjoint episodes.}
\label{tab:agentidm}
\renewcommand{\arraystretch}{1.12}
\begin{tabular*}{\linewidth}{@{\extracolsep{\fill}}llrrrr@{}}
\toprule
Readout & IDM weight & \multicolumn{2}{c}{Current state} & \multicolumn{2}{c}{Predicted state} \\
\cmidrule(lr){3-4}\cmidrule(l){5-6}
 & & MAE $\downarrow$ & $R^2\uparrow$ & MAE $\downarrow$ & $R^2\uparrow$ \\
\midrule
CLS & 0 & 44.68 & -.068 & 46.11 & -.077 \\
CLS & 3.125 & 16.78 & .703 & 26.21 & .493 \\
Spatial $8\times8$ & 0 & 45.70 & -.189 & 44.77 & -.103 \\
Spatial $8\times8$ & 3.125 & 20.87 & .680 & 25.28 & .537 \\
\bottomrule
\end{tabular*}
\end{table}

\begin{table}[!htbp]
\centering
\caption{\textbf{IDM improves planning with either readout.}
Color task success (\%); every cell reports CEM/GD. R1/R2 have 25 cases
each; Cue 1 (green) and Cue 2 (orange) have 16/34 cases, pooling ranks.
Within each readout, the IDM pair shares initialization, training budget,
and evaluation settings.}
\label{tab:idmplanning}
\renewcommand{\arraystretch}{1.12}
\setlength{\tabcolsep}{3pt}
\begin{tabular*}{\linewidth}{@{\extracolsep{\fill}}llcccc@{}}
\toprule
Readout & IDM weight & \multicolumn{2}{c}{Rank-based performance $\uparrow$} & \multicolumn{2}{c}{Cue-based performance $\uparrow$} \\
\cmidrule(lr){3-4}\cmidrule(l){5-6}
 & & R1 & R2 & Cue 1 & Cue 2 \\
\midrule
CLS & 0 & 8/12 & 0/4 & 12.5/12.5 & 0.0/5.9 \\
CLS & 3.125 & 68/52 & 8/28 & 50.0/50.0 & 32.4/35.3 \\
Spatial $8\times8$ & 0 & 12/12 & 4/4 & 12.5/6.3 & 5.9/8.8 \\
Spatial $8\times8$ & 3.125 & 80/64 & 16/52 & 56.3/62.5 & 44.1/55.9 \\
\bottomrule
\end{tabular*}
\end{table}

\FloatBarrier
\section{Spatial aggregation}
\label{app:aggregator}

Replacing CLS with spatial aggregation improves tunnel localization under
both objectives (Table~\ref{tab:readoutcontrols}). With IDM, tunnel MAE
falls from 7.42 to 3.83 pixels and $R^2$ rises from .684 to .934. This gain
concerns tunnel positions: agent-position probes are stronger with CLS+IDM
(Table~\ref{tab:agentidm}). Readout comparisons share the seed and budget,
but encoder initialization differs because constructing the aggregator
consumes random numbers.

\begin{table}[ht]
\centering
\caption{\textbf{Spatial aggregation improves tunnel localization at either IDM weight.}
Current projected states from the same four epoch-10 Color checkpoints as
Tables~\ref{tab:agentidm} and \ref{tab:idmplanning}, all trained with the same
implementation and seed 3072. MAE and $R^2$ use the same four tunnel endpoints,
indexed by cue or position; agent coordinates are evaluated separately.}
\label{tab:readoutcontrols}
\renewcommand{\arraystretch}{1.12}
\begin{tabular*}{\linewidth}{@{\extracolsep{\fill}}llrrrr@{}}
\toprule
Readout & IDM weight & Cue MAE $\downarrow$ & Cue $R^2\uparrow$ & Tunnel MAE $\downarrow$ & Tunnel $R^2\uparrow$ \\
\midrule
CLS & 0 & 31.16 & -.007 & 17.51 & -.020 \\
CLS & 3.125 & 27.44 & .307 & 7.42 & .684 \\
Spatial $8\!\times\!8$ & 0 & 28.40 & .162 & 5.03 & .898 \\
Spatial $8\!\times\!8$ & 3.125 & 27.40 & .224 & 3.83 & .934 \\
\bottomrule
\end{tabular*}
\end{table}

\subsection{Spatial readout}

Inspired by learned feature aggregation in Temporal Straightening
\citep{wang2026temporal}, we discard CLS and adaptively average-pool patch
tokens to an $8\times8$ grid. Flattening, LayerNorm, and a learned linear
map produce a 192-dimensional feature. Unlike Temporal Straightening, we
apply no temporal-curvature regularizer. Both readouts use LeWM's
Linear$(192,2048)$--BatchNorm--GELU--Linear$(2048,192)$ projector; current-state
probes act on its output.

\subsection{Why we use IDM and spatial aggregation}

IDM improves agent localization and planning, while spatial aggregation
improves tunnel localization. Together they make the cue-binding problem
easier to distinguish from broader failure: weak cue recovery persists with
CLS+IDM and with spatial features without IDM, but removing both impairs
location recovery too (Table~\ref{tab:readoutcontrols}). These controls
motivate our configuration without establishing universal necessity.
Cue-swap rollouts use only the primary spatial+IDM models; the ablations
therefore do not replicate every main-paper diagnostic.
